\documentclass[12pt,a4paper]{article}

\usepackage[T1]{fontenc}
\usepackage[utf8]{inputenc}
\usepackage[brazil]{babel}
\usepackage{lmodern}
\usepackage{geometry}
\usepackage{hyperref}
\usepackage{enumitem}
\usepackage{array}

\geometry{margin=2.5cm}
\hypersetup{
  colorlinks=true,
  linkcolor=blue,
  urlcolor=blue
}

\title{Plano de Aula: O Modelo Relacional}
\date{}

\begin{document}
\maketitle

\noindent
\textbf{Duração:} 2 horas (120 minutos)\\
\textbf{Público-alvo:} Iniciantes\\
\textbf{Recursos Necessários:} Quadro branco/negro, marcadores/giz, folhas/cartões, (opcional) projetor e computador para demonstrações.

\section{Objetivos da Aula}
\begin{itemize}[leftmargin=*,itemsep=0.25em]
  \item Apresentar os conceitos do modelo lógico relacional.
  \item Compreender relação entre Entidade (ER) e Tabela (Relacional).
  \item Identificar e diferenciar: Tabela, Coluna (Atributo), Linha (Tupla/Registro).
  \item Compreender chaves (PK/FK) e restrições básicas de integridade.
\end{itemize}

\section{Metodologia}
Aula expositiva dialogada com exemplos concretos (cadastros do dia a dia) e exercícios de transformação de descrições em tabelas. A prática é feita no quadro e em papel, reforçando estrutura e integridade.

\section{Cronograma e Conteúdo Detalhado (120 Minutos)}
\subsection*{Parte 1: Revisão e Contexto (15 minutos)}
\begin{itemize}[leftmargin=*,itemsep=0.35em]
  \item Revisar rapidamente ER: entidade, atributo, relacionamento.
  \item Pergunta-guia: ``Como isso vira armazenamento em tabelas?''
\end{itemize}

\subsection*{Parte 2: Conceitos do Modelo Relacional (30 minutos)}
\begin{itemize}[leftmargin=*,itemsep=0.35em]
  \item \textbf{Tabela (Relação) (10 min)}
  \begin{itemize}[leftmargin=*,itemsep=0.25em]
    \item Uma tabela representa um conjunto de registros do mesmo tipo.
    \item Analogia: fichário com fichas no mesmo padrão.
  \end{itemize}
  \item \textbf{Coluna e Linha (10 min)}
  \begin{itemize}[leftmargin=*,itemsep=0.25em]
    \item Coluna: o que é armazenado (Nome, Data, Valor).
    \item Linha: o item completo (um cliente, um pedido).
  \end{itemize}
  \item \textbf{Integridade e chaves (10 min)}
  \begin{itemize}[leftmargin=*,itemsep=0.25em]
    \item PK: identifica unicamente uma linha.
    \item FK: referência entre tabelas.
  \end{itemize}
\end{itemize}

\subsection*{Parte 3: Exemplo Guiado --- Clínica Veterinária (25 minutos)}
\begin{itemize}[leftmargin=*,itemsep=0.35em]
  \item Construir no quadro as tabelas CLIENTE e PET.
  \item Mostrar como FK liga PET a CLIENTE.
  \item Discutir redundância vs. separação correta.
\end{itemize}

\subsection*{Parte 4: Atividade em Grupo --- ``Desenho de Tabelas'' (35 minutos)}
\begin{itemize}[leftmargin=*,itemsep=0.35em]
  \item \textbf{Cenário:} ``Sistema de Matrícula'' (Aluno, Disciplina, Turma, Matrícula).
  \item \textbf{Passo a passo:}
  \begin{itemize}[leftmargin=*,itemsep=0.25em]
    \item Grupos de 4--5 pessoas.
    \item Definir 3--4 tabelas e seus atributos.
    \item Propor PK para cada tabela.
    \item Identificar pelo menos uma FK e explicar o relacionamento que ela representa.
  \end{itemize}
  \item \textbf{Entrega:} esquema de tabelas (nome, colunas, PK/FK).
\end{itemize}

\subsection*{Parte 5: Encerramento (15 minutos)}
\begin{itemize}[leftmargin=*,itemsep=0.35em]
  \item Revisar: tabela/coluna/linha, PK e FK.
  \item Ponte: ``Na próxima aula, vamos aprender regras para mapear ER em tabelas de forma sistemática.''
\end{itemize}

\end{document}

